Chapter 1 is being edited with Steven. Formalization has begun alongside the revised Sections 1.1--1.8, using Mathlib's standard matrices, upper-half-plane action, and Borel measures. Each section has a partial Lean companion. The chapter's status is in progress, not autoformalized.

Section 1.1 currently covers the fractional action, denominator nonvanishing, imaginary-part and difference formulas, center, kernel, and equality of kernel and center. The cocycle and central-character compatibility are deferred to Chapter 3 or later modular-forms material; the shared cocycle lemma remains available for reuse. Section 1.2 covers the elementary motions, upper-half-plane transitivity, the full stabilizer and its entry equations, the inverse position action, explicit $x,y$ coordinates, the orbit-map criterion, and existence and uniqueness of the $N,A,K$ factors as matrices. Section 1.3 proves the fractional derivative, invariant tangent coefficients via real bilinear forms, the actual real bilinear tangent form and its invariance, normalized tangent lengths, gradient duality with a real linear differential, the invariant two-point quantity $J$, and classification of positive densities satisfying the pointwise covariance law. The Laplace--Beltrami derivation, isometry commutation theorem and Green identity remain analytic formalization targets. Its manifold and measure-theoretic bridges remain open. The adjacent comparisons specify the different proof routes and list the unformalized portions, including the sphere action and scalar normalization, three-orbit theorem, quotient action and topology, angle coordinates and their classes modulo $2\pi$, and angular double covering. No statement has been weakened in the manuscript to obtain these partial implementations.

Run \texttt{python3 scripts/audit\_lean.py --chapter 1} from the project directory to build and audit the implemented portions, including Appendices A.1--A.2. The dated record is \texttt{lean/audits/chapter01-release-check-2026-11-12/record.json}; its source hashes identify the checked snapshot. This scope is partial chapter coverage, and Steven's correspondence review remains pending independently of machine checks.

Section 1.4 now covers the Cayley bijection between the actual half-plane and disk, its complex derivatives and inverse tangent metric coefficients, and the normalized point-pair invariant with its disk expression and its specialization at $i$. Rotation conjugacy, the disk/geodesic and path-distance bridges, ordered-pair orbit classification, polar area and radial integration remain open. The later definitions of $dg$ and $d\overline g$ are part of the remaining formalization target. In particular, $\theta$ is a matrix angle, the geometric rotation angle is $2\theta$, $K$ and $\overline K$ each have probability fiber measure, and discrete groups carry counting measure. The factor $|\Lambda\cap\{\pm I\}|$ in the quotient integral must be proved, not absorbed into an unspecified measure. Actual Lebesgue and quotient measures are required; a constant defined to equal the intended covolume would not formalize its calculation.

Section 1.5 now has checked algebraic ingredients: the trace discriminant and its lift independence, the interior fixed-point equation, normalization of a supplied fixed point to a rotation, exact fixed-point conditions for translations and dilations, and their point-pair displacement formulas. The dilation calculation identifies the minimizing imaginary axis for the point-pair invariant. The full boundary classification, normal-form existence from each type, actual metric infima and translation-length formulas remain open; the adjacent comparison records these distinctions.

Section 1.6 now covers the generator row operations, generation of the integral group, the two central matrix relations, integer translations, domain membership and the height bound, existence of a representative in the closed domain, and equality of two strict representatives. Full denominator estimates, centrality, boundary identifications, tiling, and the topological conditions in the closed fundamental-domain definition remain open. Section 1.7's punctured-sphere description is a prose proof, not a checked formalization.

Section 1.7 now covers the exponential cusp parameter: its definition, nonvanishing, radius, integer periods, compatibility with the action of $T^n$, exact fibers, disk preimages and limit at the puncture. These are checked complex-analytic ingredients. The stabilizer classification, definition of width as a group property, modular-neighborhood embedding, quotient homeomorphism, cross-section length, area and global topology remain open.

Section 1.8 now covers the abstract uniqueness of left Haar measure and inversion from left to right, together with the diagonal-factor argument for positive real characters: the conjugation by $k(\pi/2)$ and the consequence that a positive value equal to its inverse is $1$. Appendix A.1 contains the Borel-measure uniqueness principles and the earlier real-line lemma; Appendix A.2 contains further partial general Haar results. The group-specific Iwasawa construction, hyperbolic quotient measure, lifted-measure uniqueness, application to the actual modular function, full unimodularity and central covolumes remain open. The prose now derives invariant-area uniqueness from Haar uniqueness; the comparisons record this changed proof route.
